%% polinetwork-macros.tex
%% Definizioni per pandoc: traducono i comandi di polinetwork.cls e
%% polinetwork-commands.sty in costrutti che il reader LaTeX di pandoc capisce.
%% Viene anteposto al documento solo nella build HTML (make html-<nome>).

% Metadati: finiscono nell'header del template HTML
\newcommand{\pnTitle}[1]{\title{#1}}
\newcommand{\pnSubtitle}[1]{\subtitle{#1}}
\newcommand{\pnDate}[1]{\date{#1}}
\newcommand{\pnMakeHeader}{}
\newcommand{\pnMakeTitlePage}{}
\newcommand{\pnMakeMetadata}{}
\newcommand{\pnSetEnglish}{}
\newcommand{\pnSetItalian}{}

% pandoc scarta il gruppo che segue \noindent (es. \noindent{\small ...})
\renewcommand{\noindent}{}

% Link e testo
\newcommand{\pnHighlight}[1]{\textbf{#1}}
\newcommand{\pnWebsiteLink}{\href{https://polinetwork.org}{polinetwork.org}}
\newcommand{\pnInstagramLink}{\href{https://www.instagram.com/polinetwork_/}{@polinetwork}}
\newcommand{\pnTelegramLink}{\href{https://t.me/PoliNetwork3bot}{@PoliNetwork3bot}}
\newcommand{\pnLead}[1]{#1}
\newcommand{\pnRule}{}
